% GENERATED FILE. Edit canonical lessons, then run npm run generate:books.
% canonical-source-hash: fnv1a64:732043ab479ac7b5
% canonical-lessons: SA-C129-sitalah, SA-C129-navinah, SA-C129-pracinah, SA-C129-virah, SA-C129-caturah, SA-R129-first-pass-things-of-value-and-getting-around, SA-R129-second-pass-describing-words

\chapter{Cool, New, Ancient, Brave, Clever}
\label{ch:sa-cool-new-ancient-brave-clever}
\hlchaptermodality{129}

\begin{chapteropening}
\textbf{By the end of this chapter:} \emph{I can say cool, new, ancient, brave and clever.}

Complete the last lesson of chapter 129: I can say cool, new, ancient, brave and clever.
\end{chapteropening}

\section[\texorpdfstring{\sk{शीतलः} (śītalaḥ)}{शीतलः (śītalaḥ)}]{\sk{शीतलः} (śītalaḥ) — cool}

\label{lesson:SA-C129-sitalah}



\begin{quote}
Before the new one: say the Sanskrit for sour, then the Sanskrit for bitter.
\end{quote}

\subsection*{You'll want to know: \sk{शीतलः}}
\textbf{\sk{शीतलः}} — \emph{śītalaḥ} — \textquotedblleft{}cool\textquotedblright{}.

\begin{grammarlens}[title={What you've built}]
One more word for this chapter.
\end{grammarlens}

\subsection*{Your turn}
\begin{itemize}
\raggedright
  \item \emph{Say it:} \emph{śītalaḥ}
  \item \emph{Say it:} \emph{śītalaḥ}, once more
  \item \emph{Say it:} say \emph{śītalaḥ}
  \item \emph{Recall:} say the Sanskrit for sour, then the Sanskrit for bitter, then say \emph{śītalaḥ} again
\end{itemize}

\begin{tcolorbox}[breakable,colback=teal!4,colframe=teal!35!black,title={Before you move on}]
What does \sk{शीतलः} mean? (\textquotedblleft{}Cool\textquotedblright{}.) Say it once more.
\end{tcolorbox}
\section[\texorpdfstring{\sk{नवीनः} (navīnaḥ)}{नवीनः (navīnaḥ)}]{\sk{नवीनः} (navīnaḥ) — new}

\label{lesson:SA-C129-navinah}



\begin{quote}
Before the new one: say the Sanskrit for bitter, then the Sanskrit for cool.
\end{quote}

\subsection*{You'll want to know: \sk{नवीनः}}
\textbf{\sk{नवीनः}} — \emph{navīnaḥ} — \textquotedblleft{}new\textquotedblright{}.

\begin{grammarlens}[title={What you've built}]
One more word for this chapter.
\end{grammarlens}

\subsection*{Your turn}
\begin{itemize}
\raggedright
  \item \emph{Say it:} \emph{navīnaḥ}
  \item \emph{Say it:} \emph{navīnaḥ}, once more
  \item \emph{Say it:} say \emph{navīnaḥ}
  \item \emph{Recall:} say the Sanskrit for bitter, then the Sanskrit for cool, then say \emph{navīnaḥ} again
\end{itemize}

\begin{tcolorbox}[breakable,colback=teal!4,colframe=teal!35!black,title={Before you move on}]
What does \sk{नवीनः} mean? (\textquotedblleft{}New\textquotedblright{}.) Say it once more.
\end{tcolorbox}
\section[\texorpdfstring{\sk{प्राचीनः} (prācīnaḥ)}{प्राचीनः (prācīnaḥ)}]{\sk{प्राचीनः} (prācīnaḥ) — ancient}

\label{lesson:SA-C129-pracinah}



\begin{quote}
Before the new one: say the Sanskrit for cool, then the Sanskrit for new.
\end{quote}

\subsection*{You'll want to know: \sk{प्राचीनः}}
\textbf{\sk{प्राचीनः}} — \emph{prācīnaḥ} — \textquotedblleft{}ancient\textquotedblright{}.

\begin{grammarlens}[title={What you've built}]
One more word for this chapter.
\end{grammarlens}

\subsection*{Your turn}
\begin{itemize}
\raggedright
  \item \emph{Say it:} \emph{prācīnaḥ}
  \item \emph{Say it:} \emph{prācīnaḥ}, once more
  \item \emph{Say it:} say \emph{prācīnaḥ}
  \item \emph{Recall:} say the Sanskrit for cool, then the Sanskrit for new, then say \emph{prācīnaḥ} again
\end{itemize}

\begin{tcolorbox}[breakable,colback=teal!4,colframe=teal!35!black,title={Before you move on}]
What does \sk{प्राचीनः} mean? (\textquotedblleft{}Ancient\textquotedblright{}.) Say it once more.
\end{tcolorbox}
\section[\texorpdfstring{\sk{वीरः} (vīraḥ)}{वीरः (vīraḥ)}]{\sk{वीरः} (vīraḥ) — brave}

\label{lesson:SA-C129-virah}



\begin{quote}
Before the new one: say the Sanskrit for new, then the Sanskrit for ancient.
\end{quote}

\subsection*{You'll want to know: \sk{वीरः}}
\textbf{\sk{वीरः}} — \emph{vīraḥ} — \textquotedblleft{}brave\textquotedblright{}.

\begin{grammarlens}[title={What you've built}]
One more word for this chapter.
\end{grammarlens}

\subsection*{Your turn}
\begin{itemize}
\raggedright
  \item \emph{Say it:} \emph{vīraḥ}
  \item \emph{Say it:} \emph{vīraḥ}, once more
  \item \emph{Say it:} say \emph{vīraḥ}
  \item \emph{Recall:} say the Sanskrit for new, then the Sanskrit for ancient, then say \emph{vīraḥ} again
\end{itemize}

\begin{tcolorbox}[breakable,colback=teal!4,colframe=teal!35!black,title={Before you move on}]
What does \sk{वीरः} mean? (\textquotedblleft{}Brave\textquotedblright{}.) Say it once more.
\end{tcolorbox}
\section[\texorpdfstring{\sk{चतुरः} (caturaḥ)}{चतुरः (caturaḥ)}]{\sk{चतुरः} (caturaḥ) — clever}

\label{lesson:SA-C129-caturah}



\begin{quote}
Before the new one: say the Sanskrit for ancient, then the Sanskrit for brave.
\end{quote}

\subsection*{You'll want to know: \sk{चतुरः}}
\textbf{\sk{चतुरः}} — \emph{caturaḥ} — \textquotedblleft{}clever\textquotedblright{}.

\begin{grammarlens}[title={What you've built}]
One more word for this chapter.
\end{grammarlens}

\subsection*{Your turn}
\begin{itemize}
\raggedright
  \item \emph{Say it:} \emph{caturaḥ}
  \item \emph{Say it:} \emph{caturaḥ}, once more
  \item \emph{Say it:} say \emph{caturaḥ}
  \item \emph{Recall:} say the Sanskrit for ancient, then the Sanskrit for brave, then say \emph{caturaḥ} again
\end{itemize}

\begin{tcolorbox}[breakable,colback=teal!4,colframe=teal!35!black,title={Before you move on}]
What does \sk{चतुरः} mean? (\textquotedblleft{}Clever\textquotedblright{}.) Say it once more.
\end{tcolorbox}
\section[(dialogue)]{Things of value and getting around — a first pass}

\label{lesson:SA-R129-first-pass-things-of-value-and-getting-around}



\begin{quote}
Say the words for brave and clever.
\end{quote}

\subsection*{Your turn}
Say these aloud:

\begin{itemize}
\raggedright
  \item a coin, a post office, a letter, a book, a picture
  \item a musical instrument, a veena, a cloth, a jewel, gold
  \item silver, glass, a cart, a pot, a cot
  \item a seat, a bicycle, a vehicle, a ship, a train
  \item a cloth, a lake, a riverbank, a cuckoo, a sparrow
\end{itemize}

\begin{tcolorbox}[breakable,colback=teal!4,colframe=teal!35!black,title={Before you move on}]
A coin? (\textbf{\sk{मुद्रा}}, \emph{mudrā}.) A post office? (\textbf{\sk{पत्रालयः}}, \emph{patrālayaḥ}.) A letter? (\textbf{\sk{लेखः}}, \emph{lekhaḥ}.) A book? (\textbf{\sk{ग्रन्थः}}, \emph{granthaḥ}.) A picture? (\textbf{\sk{चित्रम्}}, \emph{citram}.) A musical instrument? (\textbf{\sk{वाद्यम्}}, \emph{vādyam}.) A veena? (\textbf{\sk{वीणा}}, \emph{vīṇā}.) A cloth? (\textbf{\sk{पट्टः}}, \emph{paṭṭaḥ}.) A jewel? (\textbf{\sk{रत्नम्}}, \emph{ratnam}.) Gold? (\textbf{\sk{सुवर्णम्}}, \emph{suvarṇam}.) Silver? (\textbf{\sk{रजतम्}}, \emph{rajatam}.) Glass? (\textbf{\sk{काचः}}, \emph{kācaḥ}.) A cart? (\textbf{\sk{शकटम्}}, \emph{śakaṭam}.) A pot? (\textbf{\sk{भाण्डम्}}, \emph{bhāṇḍam}.) A cot? (\textbf{\sk{खट्वा}}, \emph{khaṭvā}.) A seat? (\textbf{\sk{पीठम्}}, \emph{pīṭham}.) A bicycle? (\textbf{\sk{द्विचक्रिका}}, \emph{dvicakrikā}.) A vehicle? (\textbf{\sk{यानम्}}, \emph{yānam}.) A ship? (\textbf{\sk{नौयानम्}}, \emph{nauyānam}.) A train? (\textbf{\sk{धूमयानम्}}, \emph{dhūmayānam}.) A cloth? (\textbf{\sk{पटः}}, \emph{paṭaḥ}.) A lake? (\textbf{\sk{सरोवरः}}, \emph{sarovaraḥ}.) A riverbank? (\textbf{\sk{नदीतीरम्}}, \emph{nadītīram}.) A cuckoo? (\textbf{\sk{कोकिलः}}, \emph{kokilaḥ}.) A sparrow? (\textbf{\sk{चाटकः}}, \emph{cāṭakaḥ}.)
\end{tcolorbox}
\section[(dialogue)]{Describing words — a second pass}

\label{lesson:SA-R129-second-pass-describing-words}



\begin{quote}
Say the words for a mongoose and short.
\end{quote}

\subsection*{Your turn}
Say these aloud:

\begin{itemize}
\raggedright
  \item a mongoose, short, soft, clean, dirty
  \item wide, narrow, deep, wet, dry
  \item poor, rich, young, famous, safe
  \item terrible, false, pungent, sour, bitter
  \item cool, new, ancient, brave, clever
\end{itemize}

\begin{tcolorbox}[breakable,colback=teal!4,colframe=teal!35!black,title={Before you move on}]
A mongoose? (\textbf{\sk{नकुलः}}, \emph{nakulaḥ}.) Short? (\textbf{\sk{ह्रस्वः}}, \emph{hrasvaḥ}.) Soft? (\textbf{\sk{मृदुः}}, \emph{mṛduḥ}.) Clean? (\textbf{\sk{शुद्धः}}, \emph{śuddhaḥ}.) Dirty? (\textbf{\sk{मलिनः}}, \emph{malinaḥ}.) Wide? (\textbf{\sk{विस्तृतः}}, \emph{vistṛtaḥ}.) Narrow? (\textbf{\sk{संकीर्णः}}, \emph{saṅkīrṇaḥ}.) Deep? (\textbf{\sk{गभीरः}}, \emph{gabhīraḥ}.) Wet? (\textbf{\sk{आर्द्रः}}, \emph{ārdraḥ}.) Dry? (\textbf{\sk{शुष्कः}}, \emph{śuṣkaḥ}.) Poor? (\textbf{\sk{दरिद्रः}}, \emph{daridraḥ}.) Rich? (\textbf{\sk{धनिकः}}, \emph{dhanikaḥ}.) Young? (\textbf{\sk{तरुणः}}, \emph{taruṇaḥ}.) Famous? (\textbf{\sk{प्रसिद्धः}}, \emph{prasiddhaḥ}.) Safe? (\textbf{\sk{सुरक्षितः}}, \emph{surakṣitaḥ}.) Terrible? (\textbf{\sk{भयंकरः}}, \emph{bhayaṅkaraḥ}.) False? (\textbf{\sk{मिथ्या}}, \emph{mithyā}.) Pungent? (\textbf{\sk{कटुः}}, \emph{kaṭuḥ}.) Sour? (\textbf{\sk{अम्लः}}, \emph{amlaḥ}.) Bitter? (\textbf{\sk{तिक्तः}}, \emph{tiktaḥ}.) Cool? (\textbf{\sk{शीतलः}}, \emph{śītalaḥ}.) New? (\textbf{\sk{नवीनः}}, \emph{navīnaḥ}.) Ancient? (\textbf{\sk{प्राचीनः}}, \emph{prācīnaḥ}.) Brave? (\textbf{\sk{वीरः}}, \emph{vīraḥ}.) Clever? (\textbf{\sk{चतुरः}}, \emph{caturaḥ}.)
\end{tcolorbox}
